\documentclass[10pt,twocolumn]{article}
\usepackage[T1]{fontenc}
\usepackage[utf8]{inputenc}
\usepackage{lmodern}
\usepackage[margin=1.8cm,top=2.4cm,bottom=2cm,columnsep=0.6cm]{geometry}
\usepackage{fancyhdr}
\usepackage{titlesec}
\usepackage{booktabs}
\usepackage{amsmath,amssymb,amsfonts}
\usepackage{microtype}
\usepackage{xcolor}
\usepackage{mdframed}
\usepackage{parskip}
\usepackage{enumitem}
\usepackage{hyperref}

\definecolor{rulered}{RGB}{180,30,30}
\definecolor{boxgray}{RGB}{245,245,245}
\definecolor{keywordgray}{RGB}{100,100,100}

\pagestyle{fancy}
\fancyhf{}
\fancyhead[L]{\textsc{\small Claude Mag}}
\fancyhead[C]{\small\textit{Machine Learning Research Digest}}
\fancyhead[R]{\small 2026-07-29}
\fancyfoot[C]{\small\thepage}
\renewcommand{\headrulewidth}{0.8pt}

\titleformat{\section}{\normalsize\bfseries\color{rulered}}{}{0pt}{}%
  [\vspace{-4pt}\color{rulered}\rule{\columnwidth}{0.4pt}]
\titlespacing{\section}{0pt}{8pt}{4pt}

\hypersetup{colorlinks=true,urlcolor=rulered,linkcolor=black}

\begin{document}
\setlength{\parindent}{0pt}
\setlength{\parskip}{4pt}

{\small\color{keywordgray}\textbf{Keywords:}
  diffusion models \textbullet{}
  inference-time optimization \textbullet{}
  Kalman filtering \textbullet{}
  ICML 2026}\par\vspace{4pt}

{\LARGE\bfseries\raggedright
  DiFA: Inference-Time Forward-Process
  Alignment for Diffusion Models}\par\vspace{4pt}

{\large\itshape\raggedright
  Treating Each Denoising Step as a Kalman Observation Yields a
  Training-Free Temporal Consensus Wrapper with Significant FID
  Gains on CIFAR-10 and ImageNet}\par\vspace{6pt}

{\small\textbf{By} Shigui Li, Delu Zeng
  (South China University of Technology)}\par\vspace{2pt}

{\small\color{keywordgray}arXiv:2607.17972 \textbullet{}
  ICML 2026 \textbullet{} July 19, 2026}
\par\vspace{2pt}\noindent{\color{rulered}\rule{\columnwidth}{0.6pt}}\vspace{6pt}

Standard diffusion model inference treats each denoising step as an
independent numerical integration, discarding the statistical
correlation structure of predictions along the reverse trajectory.
DiFA (\textbf{Di}ffusion \textbf{F}orward-process \textbf{A}lignment),
accepted at ICML 2026, reframes inference as sequential state estimation
inspired by Kalman filtering: each prediction is an observation of the
true sample, fused with prior predictions according to noise-level
compatibility to build a forward-aligned temporal consensus. The result
is a training-free wrapper that wraps any existing sampler and yields
consistent FID improvements.

\begin{mdframed}[backgroundcolor=boxgray,linecolor=rulered,linewidth=1pt,
    innerleftmargin=6pt,innerrightmargin=6pt,
    innertopmargin=6pt,innerbottommargin=6pt]
\textbf{\small AT A GLANCE}\par\vspace{4pt}
\begin{description}[leftmargin=0pt,labelsep=4pt,itemsep=2pt,parsep=0pt,
    font=\small\bfseries]
  \item[Problem:] Inference treats denoising steps as independent
    estimators of $x_0$, ignoring forward-process statistical
    correlations across the trajectory.
  \item[Why it matters:] Exploiting the statistical structure of the
    forward process during inference is free---it requires no training
    and can improve any existing sampler.
  \item[Method:] Noise-level-weighted temporal consensus aggregates
    historical predictions; deviation guidance preserves fine-grained
    detail that averaging would smooth out.
  \item[Result:] FID 4.67$\to$3.21 on CIFAR-10 with DDIM (20 steps);
    FID 12.4$\to$8.9 on ImageNet 256$\times$256 with DDIM (50 steps).
  \item[Evidence:] Ablations confirm both the consensus and deviation
    components are necessary; uniform averaging instead of
    noise-weighted degrades FID by 0.4--1.2.
\end{description}
\end{mdframed}

\section{The Big Picture}

Diffusion models generate images by reversing a noise-injection process:
starting from Gaussian noise $x_T$, a learned denoiser iteratively
predicts the clean sample $x_0$ and steps toward it. The standard
inference framework---DDIM, DPM-Solver, EDM---casts each step as
deterministic numerical integration of the score function, treating
the model's prediction $\hat{x}_0^{(t)}$ at each step as an
independent, exact estimate.

DiFA identifies a missed opportunity: the sequence of predictions
$\{\hat{x}_0^{(T)}, \hat{x}_0^{(T-1)}, \ldots, \hat{x}_0^{(t)}\}$
is not independent---it represents a series of increasingly refined
observations of the same unknown $x_0$, each corrupted by noise
at level $\bar\alpha_t$. Kalman filtering is the optimal framework
for fusing such a sequence of noisy observations. DiFA imports this
framework into diffusion inference.

\section{Why Existing Approaches Fall Short}

The numerical integration perspective assumes each $\hat{x}_0^{(t)}$
is an unbiased, independent estimator of $x_0$. In practice:

\textbf{Early predictions are structurally aligned but coarse.}
At high noise levels $t$, the model has little to work with---predictions
are diffuse and uncertain. Yet they are statistically consistent with
the forward process, making them useful priors for later steps.

\textbf{Late predictions are fine but may drift.} At low $t$, the model
makes confident, detailed predictions---but these can deviate from the
statistical manifold of the forward process, leading to artifacts or
mode drops. Independent step treatment cannot detect or correct this drift.

\textbf{No cross-step information reuse.} Current samplers discard
$\hat{x}_0^{(\tau)}$ for $\tau > t$ once step $t$ is reached, wasting
information that has already been computed.

\section{Core Mechanism}

DiFA introduces \textbf{Forward-Process Alignment (FPA)}: at step $t$,
instead of using $\hat{x}_0^{(t)}$ directly, it computes a weighted
consensus $\tilde{x}_0^{(t)}$ across all predictions to date:

\[
  \tilde{x}_0^{(t)} = \sum_{\tau \geq t} w(\tau, t)\,\hat{x}_0^{(\tau)},
  \quad \sum_\tau w(\tau, t) = 1
\]

Weights $w(\tau, t)$ are larger for predictions at noise levels close
to $t$ (structurally compatible) and smaller for predictions made at
very different noise levels (less compatible):

\[
  w(\tau, t) \propto
    \exp\!\left(-\frac{\|\hat{x}_0^{(\tau)}-\hat{x}_0^{(t)}\|^2}{2\sigma^2}\right)
    \cdot \frac{1-\bar\alpha_\tau}{1-\bar\alpha_t}
\]

To prevent over-smoothing, DiFA adds \textbf{Deviation Guidance}:
\[
  \hat{x}_0^\text{DiFA}(t) = \tilde{x}_0^{(t)}
    + \gamma\!\left(\hat{x}_0^{(t)} - \tilde{x}_0^{(t)}\right)
\]

where $\gamma \in (0,1)$ controls the trade-off between consensus
smoothness and single-step detail.

\section{Technical Formulation}

The forward process is:
\[
  q(x_t \mid x_0) = \mathcal{N}\!\left(x_t;\,
    \sqrt{\bar\alpha_t}\,x_0,\; (1-\bar\alpha_t)I\right)
\]

Treating each $\hat{x}_0^{(\tau)}$ as an observation of $x_0$ with
additive Gaussian noise of variance proportional to $(1-\bar\alpha_\tau)$,
the minimum-variance linear estimator is:

\[
  \tilde{x}_0^{(t)} = \left(\sum_{\tau \geq t}
    \frac{1}{1-\bar\alpha_\tau}\right)^{-1}
    \sum_{\tau \geq t} \frac{\hat{x}_0^{(\tau)}}{1-\bar\alpha_\tau}
\]

DiFA approximates this by replacing the Gaussian noise assumption with
the empirical prediction discrepancy, giving the noise-and-consistency
weighted formula above. The two formulations agree when predictions are
unbiased; DiFA's version is robust to the bias introduced by imperfect
score estimation.

\section{Learning or Inference Procedure}

DiFA requires no additional training or fine-tuning. The implementation
wraps the inner loop of any existing sampler:

\begin{enumerate}[itemsep=2pt,parsep=0pt]
  \item Run the sampler's score evaluation to obtain $\hat{x}_0^{(t)}$.
  \item Append $\hat{x}_0^{(t)}$ to the history buffer.
  \item Compute $\tilde{x}_0^{(t)}$ via noise-level-weighted averaging
    over the buffer.
  \item Apply deviation guidance to obtain $\hat{x}_0^\text{DiFA}(t)$.
  \item Pass $\hat{x}_0^\text{DiFA}(t)$ to the sampler's step function
    in place of $\hat{x}_0^{(t)}$.
\end{enumerate}

Memory overhead scales as $O(T \cdot H \cdot W \cdot C)$ for the history
buffer; at 50 steps and $256 \times 256$ resolution, this is approximately
800~MB, which fits on a standard GPU.

\section{What the Guarantee Says}

Under Gaussian noise assumptions, DiFA's temporal consensus is the
minimum-variance linear estimator of $x_0$ given all observations.
This provides an optimality guarantee for the consensus itself.

For the full DiFA procedure (consensus + deviation guidance), the
optimality guarantee applies in expectation to the forward-aligned
component; the deviation term is a practical correction for the
non-Gaussian case, empirically validated but without closed-form optimality.

\section{Experimental Findings}

\textbf{CIFAR-10 unconditional (FID $\downarrow$):}

\begin{center}\small
\begin{tabular}{lccc}
\toprule
Sampler & Steps & Baseline & +DiFA \\
\midrule
DDIM & 20 & 4.67 & \textbf{3.21} \\
DPM-Solver++ & 20 & 3.44 & \textbf{2.87} \\
EDM & 35 & 2.04 & \textbf{1.78} \\
\bottomrule
\end{tabular}
\end{center}

\textbf{ImageNet 256$\times$256 class-conditional (FID $\downarrow$):}

\begin{center}\small
\begin{tabular}{lccc}
\toprule
Sampler & Steps & Baseline & +DiFA \\
\midrule
DDIM & 50 & 12.4 & \textbf{8.9} \\
DPM-Solver++ & 50 & 7.8 & \textbf{5.6} \\
\bottomrule
\end{tabular}
\end{center}

FD-DINOv2 (perceptual quality metric) and Inception Score improve
consistently across all configurations. The largest absolute FID
improvement occurs with DDIM, which makes independent step predictions
with no built-in cross-step information reuse---confirming that DiFA
fills a gap most pronounced in first-order samplers.

\section{Ablations and Interpretation}

\textbf{Pure consensus ($\gamma = 0$):} FID improves but visual
sharpness drops; generated samples appear over-smoothed compared to
reference images.

\textbf{No consensus ($\gamma = 1$):} Equivalent to baseline. The
consensus is doing the work.

\textbf{Full history vs.\ last 5 steps:} Using the full prediction
history improves FID by an additional 0.3--0.8 points, confirming that
early-trajectory predictions carry useful structural information even
when fused with weights that naturally down-weight them.

\textbf{Uniform averaging vs.\ noise-level weighting:} Replacing the
noise-level compatibility term with uniform averaging degrades FID by
0.4--1.2 points. Noise-level weighting is critical: predictions made
at very different noise levels have structurally incompatible biases
that uniform averaging cannot account for.

\textbf{Computational overhead:} Adding DiFA to DDIM (20 steps) on a
single A100 increases wall-clock inference time by 6.4\%---the
dominant cost is the buffer read and weighted average, not additional
neural function evaluations.

\section*{Reference}

Shigui Li, Delu Zeng.
\textbf{DiFA: Inference-Time Forward-Process Alignment for Diffusion Models}.
arXiv:2607.17972. ICML 2026. July 19, 2026.
\url{https://arxiv.org/abs/2607.17972}

\end{document}
